\uuid{Lf8C}
\titre{Prolongement continu et dérivées partielles en l'origine}
\niveau{L2}
\module{Analyse}
\chapitre{Fonctions de plusieurs variables}
\sousChapitre{Continuité et dérivées partielles}
\theme{Prolongement par continuité, dérivées partielles, coordonnées polaires}
\auteur{Maxime Nguyen}
\datecreate{2026-03-24}
\organisation{amscc}
\difficulte{3}

\contenu{
\texte{
  On définit $f : \mathbb{R}^2 \setminus \{(0,0)\} \to \mathbb{R}$ par
  \[
    f(x,y) = \frac{x^2}{(x^2+y^2)^{3/4}}.
  \]
}
\begin{enumerate}
  \item \question{Justifier que l'on peut prolonger $f$ en une fonction continue sur $\mathbb{R}^2$.}
    \reponse{
      On pose $x = r\cos\theta$, $y = r\sin\theta$. Alors :
      \[
        f(x,y) = \frac{r^2\cos^2\theta}{r^{3/2}} = r^{1/2}\cos^2\theta.
      \]
      Ainsi $|f(x,y)| \leq r^{1/2} \to 0$ quand $r \to 0$. Donc $\displaystyle\lim_{(x,y)\to(0,0)} f(x,y) = 0$, et on prolonge $f$ en posant $f(0,0) = 0$.
    }

  \item \question{Étudier l'existence des dérivées partielles en $(0,0)$ pour ce prolongement.}
    \reponse{
      On a $f(x,0) = \sqrt{|x|}$ pour $x \neq 0$. La fonction $x \mapsto \sqrt{|x|}$ n'est pas dérivable en $0$, donc $\dfrac{\partial f}{\partial x}(0,0)$ n'existe pas.

      En revanche, $f(0,y) = 0$ pour tout $y \in \mathbb{R}$, donc $\dfrac{\partial f}{\partial y}(0,0) = 0$.
    }
\end{enumerate}
}
